% sec-01: framing, with the positioning corrections from priority-steps.md
% applied.  Sources: potential-partners/UC-San-Diego, chunk-05 (regional
% clusters), chunk-08 (non-federal cost share), chunk-03 ($200B).
\section{What Is and Is Not Being Asked}

This is a request for a feasibility review, nothing more. No institutional
commitment, no endorsement, and no public association is sought at this stage.

Three positioning corrections are made here before anything else, because each
one would otherwise mislead a reviewer. \textbf{First}, daraxonrasib
(RMC-6236) is not first-in-human; it is investigational and already in Phase 3
evaluation. The supportable claim is the \textit{first prospective clinical
evaluation of the proposed integrated surgical and LLM advisory workflow},
subject to FDA and institutional confirmation. \textbf{Second}, the robotic
configuration must be named by manufacturer, model, cleared configuration, arm
count, instruments, software version, and intended use before any site can
assess feasibility or regulatory status; that specification is an agenda item
for the meeting, not an assertion in this document. \textbf{Third}, no drug
supply, letter of authorization, or regulatory cross-reference is in place with
the agent's developer.

\textit{Science: A New Golden Age} asks the roughly \textbf{\$200 billion}
annual federal research portfolio to invest directly in individual researchers
and to prioritise non-federal cost share and regional innovation clusters
\cite{kratsios2026sciencegoldenage,kratsiosvought2026fy2028priorities}. A San
Diego company bringing a completed regulatory package to a San Diego academic
cancer center is that pattern in its most literal form. The federal argument is
made elsewhere; what is asked here is whether the clinical concept is
supportable at all.
